\chapter{\label{ch:4-results}Results}

\minitoc

\section{Benchmarking}
To determine the effectiveness of our heuristic methods thus far,
we need to benchmark them against known methods.
In this chapter, we introduce the various benchmarking methods we use to evaluate our heuristics.
This is based off of various previous work, particularly those using ZX-calculus 
\cite{Kissinger2022cat,Kissinger_2022,sutcliffe2024proccut,Codsi2022}.
We will be using QuiZX \cite{quizx}, with features added from \cite{Kissinger2022cat},
especially the \catstate{n} decompositions.
QuiZX itself was a Rust port of PyZX, originally built in Python \cite{kissinger2020Pyzx}.

\subsection{Random Clifford+T}
Random Clifford+T circuits coming from \textit{exponentiated Pauli unitaries} 
were used in the benchmark of the BSS decomposition \cite{Kissinger2020}, 
and subsequently continued use in the benchmark of the \catstate{n} decomposition \cite{Kissinger2022cat}.
The operators in this form naturally arise in the literature, 
including Clifford+T circuit representations \cite{gosset2014}, as cited in \cite{Kissinger2020}.
We introduce their definitions, pertaining to the ZX-calculus.

\begin{definition}[Unitary]
    A unitary $U$ is a map of the form
\begin{equation}
    \tikzfig{unitary}
\end{equation} 
\end{definition}

\begin{definition}[Pauli exponential \cite{Kissinger2024}]
A \textit{Pauli exponential} is any unitary of the form
\begin{equation}
    U := e^{\pm i\theta P_1 \otimes \dots\otimes P_n}
\end{equation}
where $P_i$ are Paulis, and $\theta\in \mathbb{R}$ is a phase, and the exponential is defined as
\begin{equation}
    e^A := \sum_{k=0}^{\infty} \frac{A^k}{k!}
\end{equation}
for any map $A$, where $A^k$ is
\begin{equation}
    \cbox{\boxmap{A^k}} := \underbrace{\tikzfig{expmapa}}_k
\end{equation}
\end{definition}

\begin{proposition}[\cite{Kissinger2024}]
    Let $\vec{P} = P_1 \otimes \dots \otimes P_n$ be any Pauli string with no trivial entries, 
    so $\forall i, P_i \in \{X,Y,Z\}$.
    Then the Pauli exponential $e^{i\theta \vec{P}}$ is a ZX diagram where
\begin{equation}
    e^{\pm i\theta \vec{P}} \approx \tikzfig{pauli_exp}
\end{equation}
Here $U_j\in \{\cbox{\idwire}, \cbox{\hadamard}, \cbox{\redphase{\frac{\pi}{2}}}\}$ such that
\begin{equation}
    \tikzfig{upuj} = \cbox{\greenphase{\pi}}
\end{equation}
\end{proposition}

Choosing $\theta = (2k+1)\frac{\pi}{4}$ for $k\in \mathbb{Z}$, we can generate a random Clifford+T circuit
with a fixed depth by composing these exponentiated Pauli unitaries.
This fixes the $T$-count of the circuit, and can be used to benchmark decompositions.